\section*{The complete formal non-uniform construction (Section 4)}
The proof of Theorem 2 essentially follows Barrington's proof that the word problem of finite non-solvable groups is ALOGTIME-hard [3]. The entire proof can be found in the full version [5], where we also state a non-uniform version of Theorem 2. Since the proof for the uniform case is not difficult but tedious, in this sketch we primarily focus on the non-uniform version, which only gives us hardness via (non-uniform) $\mathrm{AC}^0$-reductions instead of DLOGTIME-reductions.
Like in Barrington's proof, we start with an ALOGTIME-machine (resp.\ $\mathrm{NC}^1$-circuit) and construct a family of so-called $G$-programs. Since we are dealing with finitely generated, but infinite groups, we have to adapt the definition of $G$-programs slightly.
Fix a finite standard generating set $\Sigma$ of $G$. A $G$-program $P$ of length $m$ and input length $n$ is a sequence of instructions $\langle i_j,b_j,c_j\rangle$ for $0\le j\le m-1$ where $i_j\in[1..n]$ and $b_j,c_j\in\Sigma$. On input of a word $x=x_1\cdots x_n\in\{0,1\}^*$, an instruction $\langle i_j,b_j,c_j\rangle$ evaluates to $b_j$ if $x_{i_j}=1$ and to $c_j$ otherwise. The evaluation of a $G$-program is the product (in the specified order) of the evaluations of its instructions, and is denoted with $P[x]\in\Sigma^*$.
Let $M$ be an ALOGTIME-machine in input normal form [47, Lemma 2.41], i.e., every computation path queries at most one input bit and $M$ halts immediately after the query. For every input size $n$, the computation tree of $M$ translates immediately into a Boolean circuit of depth $d\in O(\log n)$. Moreover, $M$ can be normalized such that this circuit is a fully balanced binary tree meaning that the gates of the circuit are indexed by the set $\{0,1\}^{\le d}$, where $\{0,1\}^{<d}$ are the inner gates (where $\varepsilon$ is the output gate, which counts here as an inner gate) and $\{0,1\}^d$ are the leaves (input gates). We can assume that all inner gates are nand-gates (where the Boolean function $\mathsf{nand}:\{0,1\}^2\to\{0,1\}$ is defined by $\mathsf{nand}(0,0)=\mathsf{nand}(0,1)=\mathsf{nand}(1,0)=1$ and $\mathsf{nand}(1,1)=0$) and each leaf is labelled by a possibly negated input variable or constant via an input mapping $q_n:\{0,1\}^d\to[1..n]\times\{0,1\}\times\{0,1\}$. The meaning of this mapping is as follows: if the input to the circuit is the bit string $x_1x_2\cdots x_n$ and $q_n(v)=\langle i,a,b\rangle$, then the input gate $v\in\{0,1\}^d$ evaluates to $a$ (resp., $b$) if $x_i=1$ (resp., $x_i=0$).
The family of circuits obtained this way can be shown to be DLOGTIME-uniform in an even stronger sense than the usual definition (see e.g.\ [47]). For the sake of a simpler description, we fix the input length $n$ and write $C$ for the $n$-input circuit of depth $d\in O(\log n)$. W.l.o.g.\ for every $x\in\{0,1\}^n$, we have $x\in L(M)$ if and only if the output gate of $C$ evaluates to $0$ on input $x$.
\begin{proof}
  For each gate $v \in \oneset{0,1}^{\leq d}$ let
  $g_v = g_{d,v}$ as in Definition 1. We construct two $G$-programs $P_v$
  and $P_v^{-1}$ (both of input length $n$) such that for every input $x \in \{0,1\}^n$
  ($x$ is taken as the input for $C$,  $P_v$, and $P_v^{-1}$) we have
  \begin{align}\label{eq:pvdef}\tag{2}
    P_v[x] &=_G \begin{cases}
      g_v &\text{if } v \text{ evaluates to } 1,\\
      1 &\text{if } v \text{ evaluates to } 0, 		
    \end{cases} 
  \end{align}
  and $P_v^{-1}[x] = P_v[x]^{-1}$ in $G$.  Notice that we have
  $g_v P_v^{-1}[x] = g_v$ if $v$ evaluates to $0$ and $g_v P_v^{-1}[x] = 1$,
  otherwise. Thus, $g_v P_v^{-1}$ is a $G$-program for the ``negation''
  of $P_v$. Moreover, by Equation (2), $P_\eps$ evaluates to $1$ on input $x$ if and only if the output gate evaluates to $0$ which by our assumption was the case if and only if  $x \in L$.
	
  The construction of the $P_v$ and $P_v^{-1}$ is straightforward:
  For an input gate $v \in \oneset{0,1}^{d}$ we simply define $P_v$ to be a $G$-program evaluating to $g_v$ or $1$~-- in which case it evaluates to which element depends on $q_n(v)$. %
  More precisely, write $g_v = a_1\cdots a_m$ with $a_i \in \Sigma$. If $q_n(v) = \langle i, a,b\rangle$ for $i \in \ccinterval{1}{n}$ and $a,b \in \oneset{0,1}$, we set 
  $P_v =  \langle i, a_1^a,a_1^b\rangle \cdots  \langle i, a_m^a,a_m^b\rangle$ and $P_v^{-1} =  \langle i, a^{-a}_m,a^{-b}_m\rangle \cdots  \langle i, a^{-a}_1,a^{-b}_1\rangle$.
  For a {\sf nand}-gate $v$ with
  inputs from $v0$ and $v1$, we define
  \begin{align*}
    P_v &= g_v [P_{v1}, P_{v0}]=g_v  P_{v1}^{-1}  P_{v0}^{-1} P_{v1} P_{v0}, \\
    P_v^{-1} &=  [P_{v0}, P_{v1}] g_v^{-1}= P_{v0}^{-1} P_{v1}^{-1}  P_{v1} P_{v1} g_v^{-1},
  \end{align*}
  where the $g_v$  and $g_v^{-1}$ represent constant
  $G$-programs evaluating to $g_v$ and $g_v^{-1}$, respectively, irrespective of the actual input
  (such constant $G$-programs consist of triples of the form $\langle 1, a,a\rangle$ for $a \in \Sigma$). These
  constant $G$-programs are defined via the commutator identities
  $g_v = \bigl[ g_{v0},\, g_{v1} \bigr]$ for $v \in \oneset{0,1}^{< d}$ in
  Definition 1.

  Clearly, by induction we have $ P_v[x]^{-1} = P_v^{-1}[x]$ in $G$ (for every input $x$). Let us
  show that Equation (2) holds: For an input gate
  $v \in \oneset{0,1}^{d}$, Equation (2) holds by definition. Now, let
  $v \in \oneset{0,1}^{< d}$. Then, by induction, we have the following equalities in $G$:
  \begin{align*}
    P_v[x] = g_v  [P_{v1}[x], P_{v0}[x]]
    &= 	\begin{cases}
      g_v  & \text{if } v0 \text{ or }  v1 \text{ evaluates to } 0, \\
      g_v [ g_{v1},\,  g_{v0} ] & \text{if } v0 \text{ and }  v1 \text{ evaluate to } 1, \\
    \end{cases}\\
    &= 	\begin{cases}
      g_v  & \text{if } v \text{ evaluates to } 1, \\
      1 & \text{if } v \text{ evaluates to } 0.		
    \end{cases}
  \end{align*}
  Note that $[ g_{v1},\,  g_{v0}] = [ g_{v0},\,  g_{v1}]^{-1} = g_v^{-1}$ for the last equality.
  Thus, $P_v$ satisfies Equation (2). For $P_v^{-1}$ the analogous
  statement can be shown with the same calculation. 
  For a leaf $v \in \{0,1\}^d$, we have $\abs{g_v} \in 2^{\mathcal{O}(d)} = n^{\mathcal{O}(1)}$
  by Condition (a) from Definition 1 (recall that $d \in \mathcal{O}(\log n)$).
  Hence, $P_v^{-1}$ and $P_v$ have polynomial
  length in $n$. Finally, also $P_\eps$ has polynomial length in $n$ (with the same
  argument as for $g_\eps$; see the remark after Definition 1).

This gives us a non-uniform $\mathrm{AC}^0$-reduction (more precisely, a projection reduction) of $L(M)$ to $\mathrm{WP}(G)$. In order to obtain a DLOGTIME-reduction, we apply essentially the same construction. However, we need to introduce some padding so that the function mapping an index $i$ encoded in binary to the $i$-th instruction of $P_\varepsilon$ can be computed in DLINTIME. In order to show the last point, we need condition (d) of the uniform SENS definition (Definition 1).
\end{proof}
